\documentclass[11pt]{article}
\usepackage[margin=1in]{geometry}
\usepackage[T1]{fontenc}
\usepackage{lmodern}
\usepackage{amsmath,booktabs,hyperref,microtype}
\hypersetup{hidelinks}
\makeatletter\let\tableinput\@@input\makeatother
% Generated; do not edit. Certified cyber-range benchmark.
\newcommand{\RangeAdaptiveRegimes}{12}
\newcommand{\RangeAdaptiveUncert}{10}
\newcommand{\RangeAdvExCoverage}{26}
\newcommand{\RangeAdvExEps}{0.10}
\newcommand{\RangeAdvExGreedy}{26}
\newcommand{\RangeAdvExK}{3}
\newcommand{\RangeAdvExOptimal}{4}
\newcommand{\RangeAttackVersion}{17.1}
\newcommand{\RangeCoverageAdaptiveCost}{15.3}
\newcommand{\RangeCoverageTypicalCost}{4.1}
\newcommand{\RangeGreedyAdaptiveCost}{17.3}
\newcommand{\RangeGreedyTypicalCost}{2.6}
\newcommand{\RangeMitigations}{44}
\newcommand{\RangeOptimalAdaptiveCost}{2.0}
\newcommand{\RangeOptimalTypicalCost}{1.6}
\newcommand{\RangeRandomAdaptiveCost}{20.3}
\newcommand{\RangeRandomTypicalCost}{10.4}
\newcommand{\RangeRegimes}{48}
\newcommand{\RangeTypicalUncert}{0}
\newcommand{\RangeUndefendedAdaptiveKOnePct}{44.9}
\newcommand{\RangeUndefendedTypicalKOnePct}{44.9}


\title{A certified cyber range: scoring automated defenders with\\
provable residual-risk guarantees on real ATT\&CK structure}
\author{}
\date{2026-09-21 (work-in-progress, version 0.1.0)}

\begin{document}
\maketitle

\begin{abstract}
Autonomous cyber-defense systems are judged almost entirely by \emph{empirical}
outcomes---attack-success rate, flags captured---which say nothing provable about
the residual risk a defender leaves behind. We describe a small \emph{cyber range}
that scores automated defenders with a \emph{provable} verdict instead. Built on the
real MITRE ATT\&CK Enterprise kill chain (v\RangeAttackVersion, \RangeMitigations{}
mitigations) and an evidence-anchored hospital-ransomware impact model, the range
returns, for any control portfolio a defender proposes, a distribution-free
\emph{guaranteed}/\emph{possible} adequacy certificate---the worst-case probability
that at least $k$ clinical services suffer simultaneous sustained outage stays below
$\varepsilon$, over all dependence structures and all parameter values in an
interval uncertainty set. We benchmark four reference automated defenders across
\RangeRegimes{} adaptive threat regimes (assumed vs.\ real CIPHER-derived
degradation, an $\varepsilon$ grid, $k=1..4$). An exact optimiser certifies at a mean
cost of \RangeOptimalMeanCost{} controls; a certificate-guided greedy defender is
close behind (\RangeGreedyMeanCost), while a certificate-blind coverage heuristic
(\RangeCoverageMeanCost) and a random policy (\RangeRandomMeanCost) pay much more,
quantifying the value of certificate feedback. The defender API is agent-agnostic, so
an agentic/LLM defender drops in unchanged---the intended next step. This is an
evaluation-environment contribution, not a new theorem or a validated real defense.
\end{abstract}

\section{The range}
A defender chooses a \emph{control portfolio}: a set of ATT\&CK mitigations
(M\#\#\#\#). Via the real \texttt{mitigates} edges, each mitigation lowers the
success probability of the techniques it covers, and the kill-chain reachability to
the ransomware impact objective composes those into a per-clinical-service sustained-
outage probability. The range's \texttt{score} is the certificate, not a sampled
outcome: worst- and best-corner evaluation over the interval uncertainty set gives a
\emph{guaranteed} verdict (adequate even at the adverse corner) and a \emph{possible}
one (adequate only at the favourable corner); the catastrophic $k$-of-$n$ event uses
the sharp closed-form aggregation bound (a min over dependence structures), so the
verdict holds whatever the true correlation between service failures. The primary
adequacy target is this catastrophic verdict, because it varies with every regime
axis. The observation (named mitigations, techniques covered, stages touched), the
action (add/remove a mitigation), and the certified score form an agent-agnostic API.

\section{Reference defenders and regimes}
We score four automated (non-LLM) defenders: an \emph{exact} minimum-cost optimiser
(budget-bounded search, feasible because adequacy is monotone in the portfolio); a
\emph{greedy} defender that adds the mitigation with the largest worst-case
reachability reduction (certificate-guided); a \emph{coverage} heuristic that maximises
techniques covered while ignoring the certificate; and a seeded \emph{random} policy.
Each is scored across \RangeRegimes{} regimes: assumed vs.\ real CIPHER-derived
degradation, adequacy levels $\varepsilon\in\{0.10,0.05,0.01\}$, and catastrophic
thresholds $k=1..4$. The reported metric is \emph{cost-to-certify}: the smallest
portfolio a defender reaches that the range certifies as adequate.

\section{Results}
Undefended, the worst-case ``at least one clinical service'' outage probability is
\RangeUndefendedKOnePct\%, far above any $\varepsilon$, so defense is required.
Table~\ref{tab:leader} ranks the defenders by cost-to-certify across the sweep: the
exact optimum averages \RangeOptimalMeanCost{} controls, the certificate-guided greedy
defender \RangeGreedyMeanCost, the certificate-blind coverage heuristic
\RangeCoverageMeanCost, and random \RangeRandomMeanCost---so \emph{using the
certificate as feedback roughly halves the control burden} relative to a
structure-only heuristic. Grounding degradation in the real CIPHER corpus rather than
the assumed prior lowers greedy's mean cost from \RangeGreedyAssumedCost{} to
\RangeGreedyCipherCost{} controls. Table~\ref{tab:gap} shows the optimality gap per
regime (assumed degradation): the greedy defender is at or within one control of the
exact optimum almost everywhere, whereas the blind heuristic's gap is larger; at the
headline regime ($\varepsilon=0.05$, $k=1$) the exact optimum is \RangeHeadOptimal{}
controls versus greedy \RangeHeadGreedy{}, coverage \RangeHeadCoverage{}, and random
\RangeHeadRandom{}. The exact optimiser is computed within the search cap in
\RangeOptimalCertPct\% of regimes (the strictest $\varepsilon$ needs larger
portfolios than the cap).

\begin{table}[t]\centering
\caption{Reference defenders by cost-to-certify across all \RangeRegimes{} regimes:
fraction of regimes certified within the search, mean controls when certified, and
the same split by assumed vs.\ real (CIPHER) degradation.}\label{tab:leader}
\begin{tabular}{@{}lrrrr@{}}\toprule
Defender & Certified (\%) & Mean controls & Assumed & CIPHER\\\midrule
\tableinput table_range_leaderboard_rows.tex
\bottomrule\end{tabular}
\end{table}

\begin{table}[t]\centering
\caption{Cost-to-certify per regime (assumed degradation): exact optimum vs.\ the
three heuristics. ``$>$4'' means the exact optimum exceeds the search cap.}\label{tab:gap}
\begin{tabular}{@{}rrrrrr@{}}\toprule
$\varepsilon$ & $k$ & Optimal & Greedy & Coverage & Random\\\midrule
\tableinput table_range_gap_rows.tex
\bottomrule\end{tabular}
\end{table}

\section{Scope and limitations}
This is a first-of-kind \emph{evaluation environment}, realistically a
work-in-progress contribution---not a new theorem (the certificate and the
closed-form $k$-of-$n$ bound are prior work) and not a validated real-world defense.
MITRE ATT\&CK and the CIPHER harm corpus are real; the hospital is synthetic, since
real hospital maps and per-incident telemetry are not public. Effectiveness and
degradation are analyst-prior intervals anchored to published figures, never incident
measurements, and the certificate verdicts are conditional worst-case frequencies
under those intervals, not statistical confidence levels. The reference defenders are
deliberately simple automated policies. The value of the range is that its score is a
\emph{provable} residual-risk bound rather than an average, and that an agentic
defender can be dropped into the same API and held to the same guarantee---the
natural next step.

\paragraph{Reproducibility.} All numbers are generated from committed tables; the
benchmark is deterministic (seeded random policy) and carries a provenance manifest
over the pinned ATT\&CK bundle and the CIPHER corpus. No LLM or network access is
used.

\end{document}
